% Options for packages loaded elsewhere
\PassOptionsToPackage{unicode}{hyperref}
\PassOptionsToPackage{hyphens}{url}
\PassOptionsToPackage{dvipsnames,svgnames,x11names}{xcolor}
\documentclass[
  11pt,
  a4paper,
]{article}
\usepackage{xcolor}
\usepackage[margin=2.3cm]{geometry}
\usepackage{amsmath,amssymb}
\setcounter{secnumdepth}{-\maxdimen} % remove section numbering
\usepackage{iftex}
\ifPDFTeX
  \usepackage[T1]{fontenc}
  \usepackage[utf8]{inputenc}
  \usepackage{textcomp} % provide euro and other symbols
\else % if luatex or xetex
  \usepackage{unicode-math} % this also loads fontspec
  \defaultfontfeatures{Scale=MatchLowercase}
  \defaultfontfeatures[\rmfamily]{Ligatures=TeX,Scale=1}
\fi
\usepackage{lmodern}
\ifPDFTeX\else
  % xetex/luatex font selection
\fi
% Use upquote if available, for straight quotes in verbatim environments
\IfFileExists{upquote.sty}{\usepackage{upquote}}{}
\IfFileExists{microtype.sty}{% use microtype if available
  \usepackage[]{microtype}
  \UseMicrotypeSet[protrusion]{basicmath} % disable protrusion for tt fonts
}{}
\makeatletter
\@ifundefined{KOMAClassName}{% if non-KOMA class
  \IfFileExists{parskip.sty}{%
    \usepackage{parskip}
  }{% else
    \setlength{\parindent}{0pt}
    \setlength{\parskip}{6pt plus 2pt minus 1pt}}
}{% if KOMA class
  \KOMAoptions{parskip=half}}
\makeatother
\usepackage{longtable,booktabs,array}
\newcounter{none} % for unnumbered tables
\usepackage{calc} % for calculating minipage widths
% Correct order of tables after \paragraph or \subparagraph
\usepackage{etoolbox}
\makeatletter
\patchcmd\longtable{\par}{\if@noskipsec\mbox{}\fi\par}{}{}
\makeatother
% Allow footnotes in longtable head/foot
\IfFileExists{footnotehyper.sty}{\usepackage{footnotehyper}}{\usepackage{footnote}}
\makesavenoteenv{longtable}
\setlength{\emergencystretch}{3em} % prevent overfull lines
\providecommand{\tightlist}{%
  \setlength{\itemsep}{0pt}\setlength{\parskip}{0pt}}
\usepackage{amsmath,amssymb}
\usepackage{booktabs}
\newcommand{\PP}{\mathbb{P}}
\newcommand{\EE}{\mathbb{E}}
\newcommand{\Var}{\operatorname{Var}}
\newcommand{\sig}{\sigma}
\usepackage{bookmark}
\IfFileExists{xurl.sty}{\usepackage{xurl}}{} % add URL line breaks if available
\urlstyle{same}
% fallback for those not using the hyperref driver hyperxmp:
\makeatletter
\@ifundefined{xmpquote}{\newcommand{\xmpquote}[1]{#1}}{}
\makeatother
\hypersetup{
  colorlinks=true,
  linkcolor={Maroon},
  filecolor={Maroon},
  citecolor={Blue},
  urlcolor={Blue},
  pdfcreator={LaTeX via pandoc}}

\author{}
\date{}

\begin{document}

\section{HW 2 --- Solutions and Marking
Scheme}\label{hw-2-solutions-and-marking-scheme}

\textbf{ID-5-02 (TO) · Mathematical Foundations of Diffusion Models ·
Monsoon 2026}

\emph{Worked solutions and marking scheme, released for student
reference.}

Total: 50 points (5 problems × 10). \textbf{Graded on correctness} ---
unlike HW 1. The per-part marks below are the marking scheme, not a
formality.

\begin{center}\rule{0.5\linewidth}{0.5pt}\end{center}

\subsection{Problem 1 --- Beta mean and variance
(10)}\label{problem-1-beta-mean-and-variance-10}

\emph{Setup:} germination rate varying across seed lots;
Beta\((\alpha,\beta)\) on \([0,1]\).

Throughout,
\(B(\alpha,\beta) = \Gamma(\alpha)\Gamma(\beta)/\Gamma(\alpha+\beta)\)
and \(\Gamma(z+1) = z\Gamma(z)\).

\textbf{(a) (1)}
\(\int_0^1 f = \frac{1}{B(\alpha,\beta)}\int_0^1 x^{\alpha-1}(1-x)^{\beta-1}dx
= \frac{B(\alpha,\beta)}{B(\alpha,\beta)} = 1\) --- immediate from the
definition of \(B\).

\textbf{(b) (3)} The integrand for \(\EE[X]\) is the Beta kernel with
\(\alpha \to \alpha+1\):

\[
\EE[X] = \frac{1}{B(\alpha,\beta)}\int_0^1 x^{\alpha}(1-x)^{\beta-1}dx
= \frac{B(\alpha+1,\beta)}{B(\alpha,\beta)}
= \frac{\Gamma(\alpha+1)}{\Gamma(\alpha)}\cdot\frac{\Gamma(\alpha+\beta)}{\Gamma(\alpha+\beta+1)}
= \frac{\alpha}{\alpha+\beta}.
\]

\emph{Marks:} 1 for recognising the shifted Beta kernel, 1 for the ratio
of \(B\)'s, 1 for applying \(\Gamma(z+1)=z\Gamma(z)\) twice to reach
\(\alpha/(\alpha+\beta)\).

\textbf{(c) (3)} Same move, shifting by two:

\[
\EE[X^2] = \frac{B(\alpha+2,\beta)}{B(\alpha,\beta)}
= \frac{\Gamma(\alpha+2)}{\Gamma(\alpha)}\cdot\frac{\Gamma(\alpha+\beta)}{\Gamma(\alpha+\beta+2)}
= \frac{\alpha(\alpha+1)}{(\alpha+\beta)(\alpha+\beta+1)} .
\]

\emph{Marks:} as in (b).

\textbf{(d) (2)}

\[
\Var(X) = \frac{\alpha(\alpha+1)}{(\alpha+\beta)(\alpha+\beta+1)} - \frac{\alpha^2}{(\alpha+\beta)^2}
= \frac{\alpha}{\alpha+\beta}\left[\frac{\alpha+1}{\alpha+\beta+1} - \frac{\alpha}{\alpha+\beta}\right]
= \boxed{\dfrac{\alpha\beta}{(\alpha+\beta)^2(\alpha+\beta+1)}}
\]

The bracket simplifies because
\((\alpha+1)(\alpha+\beta) - \alpha(\alpha+\beta+1) = \beta\).
\emph{Marks:} 1 for the subtraction set up, 1 for reaching the single
fraction. Award both if they get there by any correct route.

\textbf{(e) (1)} At \(\alpha=\beta\) the mean is
\(\alpha/(2\alpha) = 1/2\). Obvious because
\(f(x) \propto x^{\alpha-1}(1-x)^{\alpha-1}\) is then symmetric about
\(x = 1/2\). \emph{Marks:} 1, requires the symmetry reason, not just the
number.

\textbf{Common errors.} Trying to integrate \(x^\alpha(1-x)^{\beta-1}\)
by parts instead of recognising it as another Beta; forgetting the
normaliser; algebra slips in (d) --- accept any equivalent correct form.

\begin{center}\rule{0.5\linewidth}{0.5pt}\end{center}

\subsection{Problem 2 --- logit of a uniform
(10)}\label{problem-2-logit-of-a-uniform-10}

\emph{Setup:} a completely unknown probability
\(U\sim\mathrm{Uniform}(0,1)\) pushed to the log-odds scale,
\(X=\mathrm{logit}(U)\).

\textbf{(a) (1)}
\(\sig(\operatorname{logit}(u)) = \dfrac{1}{1+e^{-\log\frac{u}{1-u}}}
= \dfrac{1}{1+\frac{1-u}{u}} = \dfrac{u}{u + (1-u)} = u\). \emph{Marks:}
1.

\textbf{(b) (3)} Since \(\log\) is increasing and \(u \mapsto u/(1-u)\)
is increasing on \((0,1)\):

\[
F_X(x) = \PP\!\left(\log\tfrac{U}{1-U} \le x\right)
= \PP\!\left(\tfrac{U}{1-U} \le e^{x}\right)
= \PP\!\big(U \le \tfrac{e^{x}}{1+e^{x}}\big)
= \PP\big(U \le \sig(x)\big) = \sig(x),
\]

the last step because \(\PP(U \le c) = c\) for \(c \in (0,1)\) and
\(\sig(x) \in (0,1)\). \emph{Marks:} 1 for each rearrangement step, 1
for invoking the uniform CDF. Deduct nothing for not remarking on
monotonicity, but note it if they invert an inequality carelessly.

\textbf{(c) (3)} \(f_X(x) = F_X'(x) = \sig'(x) = \sig(x)(1-\sig(x))\).
Writing \(\sig(x) = 1/(1+e^{-x})\) and \(1-\sig(x) = e^{-x}/(1+e^{-x})\)
gives \(e^{-x}/(1+e^{-x})^2\); multiplying numerator and denominator by
\(e^{x}\) and using \(\cosh x = (e^x + e^{-x})/2\) gives
\(1/(2(\cosh x + 1))\). It integrates to \(1\) because \(F_X = \sig\)
runs from \(0\) to \(1\). \emph{Marks:} 1 for the derivative, 1 for
showing the equivalences, 1 for the normalisation argument (the CDF
argument is the cleanest; accept direct integration).

\textbf{(d) (2)}
\(f_X(-x) = \sig(-x)(1-\sig(-x)) = (1-\sig(x))\sig(x) = f_X(x)\), using
\(\sig(-x) = 1 - \sig(x)\). A symmetric density with finite mean has
mean at its centre of symmetry, so \(\EE[X] = 0\). \emph{Marks:} 1 for
the symmetry, 1 for the conclusion. They must not just assert mean
\(0\).

\textbf{(e) (1)} The \textbf{standard logistic distribution}. It was
guaranteed because \(\sig\) \emph{is} a CDF, and pushing a
\(\mathrm{Uniform}(0,1)\) through the inverse of a CDF produces exactly
the law with that CDF --- the inverse-transform sampling identity.
\emph{Marks:} 1 for the name plus the inverse-transform reason. Name
alone: 0.5, round up.

\textbf{Worth saying in class.} \(\Var(X) = \pi^2/3 \approx 3.290\);
simulation with \(4\times10^5\) draws gave mean \(+0.002\) and variance
\(3.282\). This is the same ``sampling is transport'' identity the
course uses for \(Q(U)\) --- here the target just happens to be
logistic.

\begin{center}\rule{0.5\linewidth}{0.5pt}\end{center}

\subsection{\texorpdfstring{Problem 3 --- Bernoulli to Binomial, and
\(\hat p\)
(10)}{Problem 3 --- Bernoulli to Binomial, and \textbackslash hat p (10)}}\label{problem-3-bernoulli-to-binomial-and-hat-p-10}

\emph{Setup:} A/B test on a checkout page; \(n\) independent visitors,
unknown conversion probability \(p\), \(S\) conversions.

\textbf{(a) (1)} \(\PP(x_i = x) = p^{x}(1-p)^{1-x}\) for
\(x \in \{0,1\}\). \emph{Marks:} 1. This single-expression trick is what
makes (c) painless.

\textbf{(b) (3)} (i) A specific sequence with \(k\) conversions has
probability \(p^{k}(1-p)^{n-k}\) by independence --- one factor \(p\)
per conversion and one \((1-p)\) per non-conversion, and the
\emph{order} does not change the product. (ii) There are
\(\binom{n}{k}\) such sequences (choose which of the \(n\) visitors
converted). (iii) Those sequences are disjoint events whose union is
\(\{S = k\}\), so their probabilities add. \emph{Marks:} 1 each. Part
(iii) is the one usually skipped; insist on it.

\textbf{(c) (2)}

\[
L(p) = \prod_{i=1}^n p^{x_i}(1-p)^{1-x_i} = p^{S}(1-p)^{n-S},
\qquad
\ell(p) = S\log p + (n-S)\log(1-p).
\]

\emph{Marks:} 1 each.

\textbf{(d) (3)} \(\ell'(p) = \dfrac{S}{p} - \dfrac{n-S}{1-p} = 0
\Rightarrow S(1-p) = (n-S)p \Rightarrow S = np \Rightarrow \boxed{\hat p = S/n}\).

Maximum because
\(\ell''(p) = -\dfrac{S}{p^2} - \dfrac{n-S}{(1-p)^2} < 0\) for all
\(p \in (0,1)\) (both terms negative), so \(\ell\) is strictly concave.
\emph{Marks:} 1 for the derivative, 1 for solving, 1 for the
second-derivative check. The concavity mark is not a formality --- do
not award it for ``it is obviously a max''.

Verified numerically: \(n=200\), \(S=70\); grid search over \(p\)
maximises at \(0.3500 = S/n\).

\textbf{(e) (1)} Nothing changes. The Binomial likelihood is
\(\binom{n}{S}p^{S}(1-p)^{n-S}\); the coefficient does not involve
\(p\), so it is an additive constant in \(\ell\) and vanishes on
differentiation. Same \(\hat p\). \emph{Marks:} 1, requires ``does not
depend on \(p\)''.

\textbf{Edge cases.} If \(S = 0\) or \(S = n\) the maximiser is at the
boundary (\(\hat p = 0\) or \(1\)) and \(\ell'\) never vanishes in
\((0,1)\). Worth a remark if a student raises it; not required.

\begin{center}\rule{0.5\linewidth}{0.5pt}\end{center}

\subsection{\texorpdfstring{Problem 4 --- Uniform\((a,b)\)
(10)}{Problem 4 --- Uniform(a,b) (10)}}\label{problem-4-uniformab-10}

\emph{Setup:} second-hand instrument with a lost datasheet; readings
uniform on an unknown \([a,b]\).

\textbf{(a) (3)}

\[
L(a,b) = \prod_{i=1}^n \frac{1}{b-a}\mathbf{1}\{a \le x_i \le b\}
= \frac{1}{(b-a)^n}\,\mathbf{1}\{a \le x_{(1)}\}\,\mathbf{1}\{x_{(n)} \le b\},
\]

where \(x_{(1)} = \min_i x_i\) and \(x_{(n)} = \max_i x_i\). The \(n\)
indicators collapse into two because \emph{all} observations lie in
\([a,b]\) exactly when the smallest and largest do. \emph{Marks:} 1 for
\((b-a)^{-n}\), 1 for including indicators at all, 1 for collapsing them
to the min and max. \textbf{Omitting the indicator is the central error
of this problem} --- without it the likelihood appears to increase
without bound as \(b - a \to 0\).

\textbf{(b) (2)} Two reasons, either earns the marks if argued: the
likelihood is not differentiable in \(a\) and \(b\) (it jumps to \(0\)
as they cross the data), and where it \emph{is} differentiable the
derivative never vanishes --- \((b-a)^{-n}\) is strictly decreasing in
\(b\) and strictly increasing in \(a\), so there is no interior
stationary point. The maximum sits on the boundary of the feasible
region. \emph{Marks:} 1 for identifying non-differentiability / boundary
maximum, 1 for the ``derivative never zero'' observation.

\textbf{(c) (3)} On the feasible set
\(\{a \le x_{(1)},\, b \ge x_{(n)}\}\) the likelihood is \((b-a)^{-n}\),
which is maximised by making \(b - a\) \textbf{as small as possible}.
Feasibility forces \(b - a \ge x_{(n)} - x_{(1)}\), and that lower bound
is attained at \(a = x_{(1)}\), \(b = x_{(n)}\), which is feasible.
Hence \(\hat a = \min_i x_i\), \(\hat b = \max_i x_i\). \emph{Marks:} 1
for ``minimise the width'', 1 for the feasibility bound, 1 for checking
the bound is attained.

\textbf{(d) (2)} Since \(\hat a \ge a\) and \(\hat b \le b\) always,
\(\hat b - \hat a \le b-a\) with probability \(1\) --- the estimate is
\textbf{never} too wide and is almost surely too narrow. So it is
biased, and biased \emph{downward}. \emph{Marks:} 1 for ``biased'', 1
for the correct direction with the containment reason.

\textbf{For your own reference} (not required of students): for
\(\mathrm{Uniform}(a,b)\), \(\EE[x_{(1)}] = a + \frac{b-a}{n+1}\) and
\(\EE[x_{(n)}] = b - \frac{b-a}{n+1}\), so
\(\EE[\hat b - \hat a] = \frac{n-1}{n+1}(b-a)\). Simulated at \(a=2\),
\(b=7\), \(n=10\) over \(2\times10^5\) replicates: \(\EE[\min] = 2.453\)
(theory \(2.4545\)), \(\EE[\max] = 6.547\) (theory \(6.5455\)), width
ratio \(0.8187\) (theory \(10/12 = 0.8182\)). The unbiased correction
multiplies the width by \((n+1)/(n-1)\).

\begin{center}\rule{0.5\linewidth}{0.5pt}\end{center}

\subsection{Problem 5 --- cross-entropy from maximum likelihood
(10)}\label{problem-5-cross-entropy-from-maximum-likelihood-10}

\emph{Setup:} postal digit sorter, \(K=10\) classes, logits
\(z_i=f_\theta(x_i)\in\mathbb{R}^K\),
\(q_{ik}=e^{z_{ik}}/\sum_m e^{z_{im}}\), labels \(y_i\) one-hot.

\textbf{(a) (2)} Positivity: \(e^{z_{ik}}>0\) for every real \(z_{ik}\),
and the denominator is a sum of positive terms, so \(q_{ik}>0\).
Normalisation:

\[
\sum_{k=1}^{K} q_{ik}
= \frac{\sum_{k=1}^{K} e^{z_{ik}}}{\sum_{m=1}^{K} e^{z_{im}}} = 1 ,
\]

the numerator and denominator being the same sum. Shift-invariance:
replacing \(z_{ik}\) by \(z_{ik}+c\),

\[
\frac{e^{z_{ik}+c}}{\sum_m e^{z_{im}+c}}
= \frac{e^{c}\,e^{z_{ik}}}{e^{c}\sum_m e^{z_{im}}}
= q_{ik} .
\]

So softmax is constant along the all-ones direction: only
\textbf{\(K-1\)} of the \(K\) logits are genuinely free. The map
\(\mathbb{R}^K \to\) (the \((K-1)\)-simplex) is onto but not injective,
which is why implementations may pin \(z_K = 0\) without losing
anything.

\emph{Marks:} 1 for positivity \textbf{and} normalisation together, 1
for shift-invariance with the \(K-1\) conclusion. Accept ``the extra
degree of freedom is unidentifiable''.

\textbf{(b) (2)} \(y_i \mid x_i \sim \mathrm{Categorical}(q_i)\) ---
equivalently \(\mathrm{Multinomial}(1, q_i)\) --- with pmf

\[
\PP(y_i \mid x_i, \theta) = \prod_{k=1}^{K} q_{ik}^{\,y_{ik}} .
\]

Because \(y_i\) is one-hot, every factor with \(y_{ik}=0\) is
\(q_{ik}^{0}=1\), so the product collapses to \(q_{i c_i}\), the
predicted probability of the true class \(c_i\). This is the exact
\(K\)-class analogue of \(q^{y}(1-q)^{1-y}\).

\emph{Marks:} 1 for naming Categorical/Multinoulli with \(q_i\) as
parameter, 1 for the single-expression pmf. It must be
\textbf{conditional on \(x_i\)}. Writing \(q_{i c_i}\) directly is fine
and earns full marks if the one-hot product is also given or clearly
implied.

\textbf{(c) (3)}

\[
L(\theta) = \prod_{i=1}^{n}\prod_{k=1}^{K} q_{ik}^{\,y_{ik}},
\qquad
\ell(\theta) = \sum_{i=1}^{n}\sum_{k=1}^{K} y_{ik}\log q_{ik},
\qquad
-\tfrac{1}{n}\ell(\theta) = \mathrm{CE}(\theta).
\]

So \textbf{minimising cross-entropy is exactly maximising the
likelihood} of the observed labels under the assumed Categorical model;
\(-1/n\) is a monotone rescaling, so it moves the objective's value but
not its argmax.

The three assumptions:

\begin{enumerate}
\def\labelenumi{\arabic{enumi}.}
\tightlist
\item
  \textbf{Form} --- each label is Categorical, with its class
  probabilities given by the model output
  \(\mathrm{softmax}(f_\theta(x_i))\). (Two things bundled: that the
  label is Categorical at all, and that the softmax link is the right
  parametrisation.)
\item
  \textbf{Across examples} --- the labels are conditionally independent
  given the inputs, which is what licenses the product over \(i\).
\item
  \textbf{The inputs} --- \(x_i\) are treated as fixed/given, not
  modelled. Nothing here assigns a distribution to \(x\), so this is a
  \emph{conditional} (discriminative) likelihood, not a joint model of
  \((x,y)\).
\end{enumerate}

\emph{Marks:} 1 for the likelihood and log, 1 for landing on
\(\mathrm{CE}\) with the argmax remark, 1 for the assumptions --- award
the full point only if all three appear. Assumption 3 is the one almost
everyone omits and the one worth drawing out in class.

\textbf{(d) (2)} Write the log-softmax in the form that makes the
derivative trivial:

\[
\log q_k = z_k - \log\!\sum_{m} e^{z_m}
\;\Longrightarrow\;
\frac{\partial \log q_k}{\partial z_j}
= \delta_{kj} - \frac{e^{z_j}}{\sum_m e^{z_m}}
= \delta_{kj} - q_j .
\]

Then

\[
\frac{\partial \ell}{\partial z_j}
= -\sum_{k} y_k\big(\delta_{kj} - q_j\big)
= -y_j + q_j \underbrace{\sum_k y_k}_{=\,1}
= q_j - y_j .
\]

The constraint \(\sum_k y_k = 1\) is used at exactly one place: pulling
\(q_j\) out of the sum and collapsing what remains. Without one-hot (or
at least sum-to-one) labels the second term would carry a factor
\(\sum_k y_k \ne 1\).

\emph{Marks:} 1 for \(\partial \log q_k/\partial z_j = \delta_{kj}-q_j\)
(however obtained --- quotient rule is fine but slower), 1 for the
collapse to \(q_j - y_j\) \textbf{with} the \(\sum_k y_k = 1\) step
identified. Do not award the second mark if the student writes
\(q_j - y_j\) without ever using the constraint.

\textbf{(e) (1)} With \(K=2\),

\[
q_2 = \frac{e^{z_2}}{e^{z_1}+e^{z_2}}
    = \frac{1}{1+e^{-(z_2-z_1)}} = \sig(z),
\qquad z := z_2 - z_1,
\]

after dividing numerator and denominator by \(e^{z_2}\); and
\(q_1 = 1 - \sig(z)\). Writing \(y \in \{0,1\}\) for the indicator of
class 2, the one-hot label is \((1-y, y)\) and

\[
-\big[(1-y)\log q_1 + y\log q_2\big]
= -\big[y\log\sig(z) + (1-y)\log(1-\sig(z))\big],
\]

which is BCE. Two logits were one too many precisely by the
shift-invariance of (a): only the \emph{difference} \(z_2 - z_1\) is
identifiable.

\emph{Marks:} 1 for the reduction with the shift-invariance link.
Getting \(q_2 = \sig(z_2-z_1)\) but not connecting it to (a) earns the
mark only if the BCE collapse is also shown.

\textbf{Why this is worth the time.} ``Predicted minus observed'', now
in every coordinate at once, is why softmax classifiers train as easily
as they do, and the same structure reappears throughout the course: the
two-point posterior \(\PP(x_0 = +1 \mid x_\sigma) = \sig(2x/\sigma^2)\)
is the \(K=2\) case, and the denoiser's loss has the same
predicted-minus-observed shape. The binary case is derived in the
``binary cross-entropy is a likelihood'' aside in tutorial 02 --- this
problem is the general statement that aside is a special case of.

\begin{center}\rule{0.5\linewidth}{0.5pt}\end{center}

\subsection{Marking summary}\label{marking-summary}

{\def\LTcaptype{none} % do not increment counter
\begin{longtable}[]{@{}llr@{}}
\toprule\noalign{}
Problem & Topic & Points \\
\midrule\noalign{}
\endhead
\bottomrule\noalign{}
\endlastfoot
1 & Beta: mean and variance by integration & 10 \\
2 & Inverse sigmoid of a uniform → logistic & 10 \\
3 & Bernoulli → Binomial; MLE of \(p\) & 10 \\
4 & Uniform\((a,b)\): MLE from the support & 10 \\
5 & Softmax cross-entropy as maximum likelihood & 10 \\
& \textbf{Total} & \textbf{50} \\
\end{longtable}
}

\subsubsection{Errors to expect}\label{errors-to-expect}

\begin{enumerate}
\def\labelenumi{\arabic{enumi}.}
\tightlist
\item
  \textbf{Omitting the indicator in 4(a)} --- the defining error of that
  problem.
\item
  Treating 4 as a calculus problem and reporting nonsense stationary
  points.
\item
  In 3(b), skipping the disjointness step and just asserting the
  coefficient.
\item
  In 5(c), missing that the inputs are unmodelled --- the discriminative
  / generative distinction.
\item
  In 5(d), producing \(q_j - y_j\) without ever invoking
  \(\sum_k y_k = 1\), or grinding through the quotient rule instead of
  using \(\log q_k = z_k - \log\sum_m e^{z_m}\).
\item
  In 5(a), verifying normalisation but skipping shift-invariance, and so
  missing that only \(K-1\) logits are free.
\item
  In 1, attempting integration by parts rather than recognising the Beta
  kernel.
\item
  Asserting \(\EE[X] = 0\) in 2(d) without the symmetry argument.
\end{enumerate}

\end{document}
